\begin{frame} 
\frametitle{Tipologie di messaggi ICMP} 
Le linee guida RFC 792 (ICMPv4) e RFC 4443 (ICMPv6) definiscono i messaggi presenti in ICMP.  
E li classifica o come messaggi di errore, che segnalano problemi nella comunicazione di rete, o come messaggi informativi, 
che vengono utilizzati per scopi diagnostici e di controllo. 
\end{frame} 
\begin{frame} 
\frametitle{Destination Unreachable} 
Viene inviato quando il pacchetto non può essere recapitato alla destinazione specificata nell'header IP. 
Di solito perchè il percorso definito non può essere seguito.  
Il codice impostato nel messaggio indicherà il motivo dell'invio.  
\begin{bytefield}[bitwidth=1.1em]{32} 
	%\bitbox{8}{0} & \bitbox{8}{1} & \bitbox{8}{2} & \bitbox{8}{3} \\
	\bitheader{0-31} \\
	\bitbox{8}{Type=3 (1 byte)} & \bitbox{8}{Code=0-5 (1 byte)} & \bitbox{16}{Checksum (2 byte)} \\
	\bitbox{32}{Unused (4 byte)} \\
	\bitbox{32}{Internet Header + 64 bits of Original Datagram ($\geq$ 21 byte)} 
\end{bytefield}  
%\captionof{figure}{Struttura del pacchetto \textbf{ICMPv4 Destination Unreachable}} 
%Il campo \textit{Internet Header}: viene utilizzato dall'host per accoppiare il messaggio di errore al processo appropriato. 
\begin{bytefield}[bitwidth=1.1em]{32} 
	%\bitbox{8}{0} & \bitbox{8}{1} & \bitbox{8}{2} & \bitbox{8}{3} \\
	\bitheader{0-31} \\
	\bitbox{8} {Type=1 (1 byte)} & \bitbox{8}{Code=0-6 (1 byte)} & \bitbox{16}{Checksum (2 byte)} \\
	\bitbox{32} {Unused (4 byte)} \\
	\bitbox{32} {As much of invoking packet as possible without} \\
	\bitbox{32} {the ICMPv6 packet exceeding the minimum IPv6 MTU ($\geq$ 0 byte)} 
\end{bytefield} 
%\captionof{figure}{Struttura del pacchetto \textbf{ICMPv6 Destination Unreachable}} 
%Il campo \textit{Invoking Packet} indica quanta parte del pacchetto (che ha attivato l'errore ICMPv6) debba essere inclusa. 
%Il tutto senza eccedere il \textit{IPv6 MTU} il cui valore di default equivale a 1280 bytes. 
\end{frame} 

\begin{frame} 
\frametitle{Time Exceeded}  
Questa tipologia di messaggio viene usata quando il gateway che elabora un pacchetto trova che il suo TTL (tempo di vita) è zero. 
In questi casi il gateway dovrà scartare il datagramma e notificare l'host sorgente della cosa. 
\begin{bytefield}[bitwidth=1.1em]{32} 
        %\bitbox{8}{0} & \bitbox{8}{1} & \bitbox{8}{2} & \bitbox{8}{3} \\
        \bitheader{0-31} \\
        \bitbox{8}{Type=11 (1 byte)} & \bitbox{8}{Code=0-1 (1 byte)} & \bitbox{16}{Checksum (2 byte)} \\
        \bitbox{32}{Unused (4 byte)} \\
        \bitbox{32}{Internet Header + 64 bits of Original Datagram ($\geq$ 21 byte)} 
\end{bytefield} 
\captionof{figure}{Struttura del pacchetto \textbf{ICMPv4 Time Exceeded}}
\begin{bytefield}[bitwidth=1.1em]{32} 
        %\bitbox{8}{0} & \bitbox{8}{1} & \bitbox{8}{2} & \bitbox{8}{3} \\
        \bitheader{0-31} \\
        \bitbox{8} {Type=3 (1 byte)} & \bitbox{8}{Code=0-1 (1 byte)} & \bitbox{16}{Checksum (2 byte)} \\
        \bitbox{32} {Unused (4 byte)} \\
        \bitbox{32} {As much of invoking packet as possible without} \\
        \bitbox{32} {the ICMPv6 packet exceeding the minimum IPv6 MTU ($\geq$ 0 byte)} 
\end{bytefield}  
\captionof{figure}{Struttura del pacchetto \textbf{ICMPv6 Time Exceeded}}
\end{frame} 

\begin{frame} 
\frametitle{Parameter Problem}  
Viene usata quando il gateway che elabora un pacchetto  trova un problema con i parametri dell'intestazione 
in modo tale da non poter completare l'elaborazione del datagramma. 
In questo caso dovrà scartare il datagramma e notificare la cosa all'host indicando il tipo e la posizione del problema. 
Il messaggio viene inviato solo se l'errore ha causato lo scarto del pacchetto. 
\begin{bytefield}[bitwidth=1.1em]{32} 
        %\bitbox{8}{0} & \bitbox{8}{1} & \bitbox{8}{2} & \bitbox{8}{3} \\
        \bitheader{0-31} \\
        \bitbox{8}{Type=12 (1 byte)} & \bitbox{8}{Code=0 (1 byte)} & \bitbox{16}{Checksum (2 byte)} \\
        \bitbox{8}{Pointer (1 byte)} & \bitbox{24}{Unused (3 byte)} \\
        \bitbox{32}{Internet Header + 64 bits of Original Datagram ($\geq$ 21 byte)} 
\end{bytefield}  
\captionof{figure}{Struttura del pacchetto \textbf{ICMPv4 Parameter Problem}}
\begin{bytefield}[bitwidth=1.1em]{32} 
        \bitheader{0-31} \\
        \bitbox{8} {Type=4 (1 byte)} & \bitbox{8}{Code=0-2 (1 byte)} & \bitbox{16}{Checksum (2 byte)} \\
        \bitbox{32} {Pointer (4 byte)} \\
        \bitbox{32} {As much of invoking packet as possible without} \\
        \bitbox{32} {the ICMPv6 packet exceeding the minimum IPv6 MTU ($\geq$ 0 byte)} 
\end{bytefield} 
\captionof{figure}{Struttura del pacchetto \textbf{ICMPv6 Parameter Problem}}
\end{frame} 

\begin{frame} 
\frametitle{Source Quench}  
Questa tipologia viene usata quando il gateway vuole richiedere di ridurre la velocità di invio dei pacchetti. 
Questo perchè il gateway ha scartato un pacchetto ma non a causa di un errore.  
Al suo ricevimento, l'host sorgente dovrà ridurre la velocità sino a quando non 
riceverà più messaggi del tipo Source Quench dal gateway. 
\begin{bytefield}[bitwidth=1.1em]{32} 
        %\bitbox{8}{0} & \bitbox{8}{1} & \bitbox{8}{2} & \bitbox{8}{3} \\
        \bitheader{0-31} \\
        \bitbox{8}{Type=4 (1 byte)} & \bitbox{8}{Code=0 (1 byte)} & \bitbox{16}{Checksum (2 byte)} \\
        \bitbox{32}{Unused (4 byte)} \\
        \bitbox{32}{Internet Header + 64 bits of Original Datagram ($\geq$ 21 byte)} 
\end{bytefield} 
\captionof{figure}{Struttura del pacchetto \textbf{ICMPv4 Source Quench}} 
\end{frame} 

\begin{frame} 
\frametitle{Redirect}  
Indica un messaggio di reindirizzamento a un host.  
Il gateway manda questo tipo di messaggio se, dopo aver controllato la sua tabella di routing, 
trova che esiste un gateway migliore che si trova sulla sua stessa rete. 
Questo secondo gateway rappresenterà un percorso migliore per la destinazione.  
Il codice impostato nel messaggio indicherà uno dei motivi presenti nella Tabella \ref{tabella:caso:Redirect}. 
\begin{bytefield}[bitwidth=1.1em]{32} 
        %\bitbox{8}{0} & \bitbox{8}{1} & \bitbox{8}{2} & \bitbox{8}{3} \\
        \bitheader{0-31} \\
        \bitbox{8}{Type=5 (1 byte)} & \bitbox{8}{Code=0-3 (1 byte)} & \bitbox{16}{Checksum (2 byte)} \\
        \bitbox{32}{Gateway Internet Address (4 byte)} \\
        \bitbox{32}{Internet Header + 64 bits of Original Datagram ($\geq$ 21B)} 
\end{bytefield}
\captionof{figure}{Struttura del pacchetto \textbf{ICMPv4 Redirect}}
\end{frame} 

\begin{frame} 
\frametitle{Echo Request/Echo Reply}  
Un messaggio \textit{Echo}, viene usato per ricevere indietro una risposta da un host. 
Si inviano dei dati tramite una Echo Request, e questi'ultimi dovranno essere restituiti in un messaggio di risposta integralmente e senza modifiche. 
\begin{bytefield}[bitwidth=1.1em]{32} 
        %\bitbox{8}{0} & \bitbox{8}{1} & \bitbox{8}{2} & \bitbox{8}{3} \\
        \bitheader{0-31} \\
        \bitbox{8}{Type=8 (1 byte)} & \bitbox{8}{Code=0 (1 byte)} & \bitbox{16}{Checksum (2 byte)} \\
        \bitbox{16}{Identifier (2 byte)} && \bitbox{16}{Sequence Number (2 byte)} \\
        \bitbox{32}{Data ... ($\geq$ 0 byte)} 
\end{bytefield}  
\captionof{figure}{Struttura del pacchetto \textbf{ICMPv4 Echo}}
\begin{bytefield}[bitwidth=1.1em]{32} 
        %\bitbox{8}{0} & \bitbox{8}{1} & \bitbox{8}{2} & \bitbox{8}{3} \\
        \bitheader{0-31} \\
        \bitbox{8}{Type=128 (1 byte)} & \bitbox{8}{Code=0 (1 byte)} & \bitbox{16}{Checksum (2 byte)} \\
        \bitbox{16}{Identifier (2 byte)} && \bitbox{16}{Sequence Number (2 byte)} \\
        \bitbox{32}{Data ... ($\geq$ 0B)} 
\end{bytefield} 
\captionof{figure}{Struttura del pacchetto \textbf{ICMPv6 Echo}}
Il mittente non ha alcuna limitazione sulla quantità di dati inseribili  nel campo del payload. 
Tuttavia una limitazione potrà essere data dalla massima capacità di trasporto dei collegamenti. 
Nel caso la dimensione del messaggio la superasse; il pacchetto dovrà essere frammentato per poter essere spedito. 
\end{frame} 

\begin{frame} 
\frametitle{Timestamp Request /Timestamp Reply}  
Viene usato per ricevere indietro una risposta da un host. 
I dati ricevuti nel messaggio di richiesta, vengono restituiti in quello di risposta insieme a dei timestamp aggiuntivi. 
Il timestamp è pari a 32 bit e indica i millisecondi che sono passati dalla mezzanotte UT.
\begin{bytefield}[bitwidth=1.1em]{32} 
        %\bitbox{8}{0} & \bitbox{8}{1} & \bitbox{8}{2} & \bitbox{8}{3} \\
        \bitheader{0-31} \\
        \bitbox{8}{Type=13 (1 byte)} & \bitbox{8}{Code=0 (1 byte)} & \bitbox{16}{Checksum (2 byte)} \\
        \bitbox{16}{Identifier (2 byte)} && \bitbox{16}{Sequence Number (2 byte)} \\ 
        \bitbox{32}{Originate Timestamp (4 byte)} \\
        \bitbox{32}{Receive Timestamp (4 byte)} \\
        \bitbox{32}{Transmit Timestamp (4 byte)} \\
        \bitbox{32}{Data ... ($\geq$ 0 byte)} 
\end{bytefield}
\captionof{figure}{Struttura del pacchetto \textbf{ICMPv4 Timestamp}} 
Mentre i campi relativi ai timestamp indicheranno rispettivamente: 
il tempo in cui il mittente ha toccato il messaggio per l'ultima volta prima di inviarlo, 
il tempo in cui il destinatario ha toccato per la prima volta il messaggio (alla ricezione) e 
il tempo in cui il destinatario ha toccato il messaggio per l'ultima volta prima di inviarlo.  
\end{frame} 

\begin{frame} 
\frametitle{Information Request /Information Reply}  
La tipologia \textit{Information} viene usata per consentire di scoprire il numero della rete in cui un host si trova. 
Serve quindi per capire se si trova nella stesse rete dell'host che risponde. 
Sebbene il messaggio può essere inviato con la destinazione nell'intestazione IP pari a zero (ciò significa "questa" rete); 
l'intestazione IP presente nel messaggio di risposta dovrà essere inviata con gli indirizzi IP completamente specificati.
\begin{bytefield}[bitwidth=1.1em]{32} 
        %\bitbox{8}{0} & \bitbox{8}{1} & \bitbox{8}{2} & \bitbox{8}{3} \\
        \bitheader{0-31} \\
        \bitbox{8}{Type=15 (1 byte)} & \bitbox{8}{Code=0 (1 byte)} & \bitbox{16}{Checksum (2 byte)} \\
        \bitbox{16}{Identifier (2 byte)} && \bitbox{16}{Sequence Number (2 byte)} 
\end{bytefield} 
\captionof{figure}{Struttura del pacchetto \textbf{ICMPv4 Information}}
\end{frame} 

\begin{frame} 
\frametitle{Packet Too Big}  
Un messaggio \textit{Packet Too Big} viene generato da un router in risposta a un pacchetto che 
non può inoltrare perché è più grande dell'MTU del collegamento in uscita.
\begin{bytefield}[bitwidth=1.1em]{32} 
        %\bitbox{8}{0} & \bitbox{8}{1} & \bitbox{8}{2} & \bitbox{8}{3} \\
        \bitheader{0-31} \\
        \bitbox{8} {Type=2 (1 byte)} & \bitbox{8}{Code=0 (1 byte)} & \bitbox{16}{Checksum (2 byte)} \\
        \bitbox{32} {MTU (4 byte)} \\
        \bitbox{32} {As much of invoking packet as possible without} \\
        \bitbox{32} {the ICMPv6 packet exceeding the minimum IPv6 MTU ($\geq$ 0 byte)} 
    \end{bytefield} 
    \captionof{figure}{Struttura del pacchetto \textbf{ICMPv6 Packet Too Big}}

\end{frame} 
\begin{frame} 
\frametitle{}  

\end{frame} 
\begin{frame} 
\frametitle{}  

\end{frame} 
\begin{frame} 
\frametitle{}  

\end{frame} 
\begin{frame} 
\frametitle{}  

\end{frame} 